\begin{claim}{D.57}{Proved}{Full lifts of the noncyclic torus two-groups}
Let $Q$ be a subgroup of the full
once-punctured-torus loop-outer group,
isomorphic to $C_2\times C_2$ or to
$D_8$ of order eight.  Put
\[
 E_Q=\{(\phi,q)\in\operatorname{Aut}(F)\times Q:
                        [\phi]=q\},\qquad F=\widehat F(x,y).
                                                               \tag{D.57a}
\]
The projection $E_Q\to Q$ has a homomorphic
section with exactly its specified outer labels.
For $Q=D_8$ there is exactly one
$F$-conjugacy class of such sections.

For $Q=C_2\times C_2$, let $z$ be its
positive involution and let $r$ be either
of its two negative involutions.  There
are exactly three section classes when
\[
                         A_2(r)\bmod2=I,                \tag{D.57b}
\]
and exactly one otherwise.  This choice
does not depend on which negative
involution is called $r$.
For the image $H$ of every full section,
\[
                  N_{E_Q}(H)=C_F(H)\times H,
                  \qquad H^1(C_F(H),\F_2)=0.            \tag{D.57c}
\]
Together with D.27 and D.46, this supplies
actual lifts for every finite two-subgroup
in the list of D.50.  It is a statement
about labelled automorphism lifts, not
a geometric conjugacy classification.
\end{claim}
\begin{proof}
\emph{The invariant prime-section classes for a four-group.}
Suppose first that $Q=C_2\times C_2$.
By D.50 and D.55 its faithful dyadic
representation has exactly one nonidentity
positive element $z$, with $A_2(z)=-I$;
the two other involutions are negative.
In the inverse image $E_z$ of
$\langle z\rangle$, D.46 supplies exactly
three $F$-conjugacy classes of sections
$T$ of order two, each with
\[
                         H^1(C_F(T),\F_2)=0.             \tag{D.57d}
\]
As in (D.56d), their elementary abelian
coordinates are three distinct points
of the affine plane $\F_2^2$.

Conjugation by an element of $E_Q$
with label $r$ induces an affine map
\[
                           v\longmapsto Mv+d,
                    \qquad M=A_2(r)\bmod2.              \tag{D.57e}
\]
Its square is the identity on the plane:
the square of the conjugating element
lies in $F$, and inner automorphisms
act trivially on the elementary abelian
quotient of $E_z$.
Its three-point set of section classes
is invariant.

If $M=I$, a nonzero translation would
have only two-element orbits, so could
not preserve a three-point set.
Thus $d=0$, and all three classes are
fixed. If $M\ne I$, it is the unique
two-dimensional unipotent type over
$\F_2$, with
\[
          \ker(I+M)=\operatorname{im}(I+M),
                         \quad\dim\ker(I+M)=1.
\]
The square identity gives $(I+M)d=0$,
so (D.57e) has two fixed points and
one orbit of size two.
An invariant three-point subset must
contain that orbit and one fixed
point. Thus exactly one section
class is fixed in this case.
Since $A_2(zr)=-A_2(r)$, the reductions
of the two negative involutions are
equal, proving the independence in
(D.57b).

Choose any fixed class and one of
its section groups $T$.
Its normalizer $N=N_{E_Q}(T)$ maps
onto $Q$: adjust a conjugator with
label $r$ by an element of $F$ to
normalize $T$, and use $T$ itself
for the label $z$.
Its kernel is $C=C_F(T)$, since an
element of $F$ normalizing the
order-two group $T$ centralizes it.
Moreover $T$ is central in $N$.
The quotient $J=N/T$ therefore has
the exact sequence
\[
                  1\longrightarrow C\longrightarrow J
                     \longrightarrow Q/\langle z\rangle
                     \longrightarrow1,
                  \qquad Q/\langle z\rangle=C_2.        \tag{D.57f}
\]
The inclusion of $C$ in this quotient
is injective, because $C\cap T=1$.

\emph{The invariant order-four class for a dihedral group.}
Now suppose that $Q=D_8$ and let
$R\cong C_4$ be its rotation subgroup.
By D.53 the inverse image $E_R$ has
exactly one $F$-conjugacy class of
labelled order-four sections $B$,
and
\[
                         H^1(C_F(B),\F_2)=0.             \tag{D.57g}
\]
Every element of $E_Q$ preserves this
class, since $R$ is normal in $Q$.
The normalizer $N=N_{E_Q}(B)$ thus
maps onto $Q$, by the same adjustment
of a conjugator by an element of $F$.
Its kernel is $C=C_F(B)$: an element
of $F$ normalizing $B$ fixes the label
of every element of $B$, and therefore
centralizes $B$.
The subgroup $B$ is normal in $N$,
and the quotient $J=N/B$ has
\[
                  1\longrightarrow C\longrightarrow J
                        \longrightarrow Q/R\longrightarrow1,
                              \qquad Q/R=C_2.           \tag{D.57h}
\]
Again $C$ embeds, because $C\cap B=1$.
Conjugation on $B$ in this construction
is the actual conjugation on its
specified rotation labels in $D_8$.

\emph{Torsion and exact finite preimages in the two constructed quotients.}
In each of (D.57f) and (D.57h), the
group $C$ is a closed subgroup of
$F$, and hence is torsion-free with
$\operatorname{cd}_p(C)\leq1$ at every
prime. Equations (D.57d) and (D.57g)
give $C^{(2)}=1$.
Every map from $J$ to a finite two-group
therefore kills $C$, and the displayed
nontrivial quotient gives
\[
                              J^{(2)}=C_2.              \tag{D.57i}
\]
The same full cohomological-dimension
argument used in D.27 and D.46 proves
that $J$ has torsion. Indeed, a
torsion-free $J$ would be projective,
since $C$ is open and has the stated
cohomological dimensions. Its maximal
pro-two quotient would then be a free
pro-two group, contrary to (D.57i).
Every finite subgroup of $J$ injects
into $J/C=C_2$. Consequently $J$ has
an involution $j$ outside $C$.

Take the full inverse image $H$ of
$\langle j\rangle$ in $N$.
In the four-group construction its
kernel is $T$, so $|H|=4$; in the
dihedral construction its kernel is
$B$, so $|H|=8$.
In either case $H\cap F=1$.
To check this in the actual maps,
an element of $H\cap F$ has image
in $\langle j\rangle\cap C=1$ in
$J$, and then lies in $T\cap F=1$
or in $B\cap F=1$ respectively.
Thus the projection $H\to Q$ is
injective. Its order equals $|Q|$,
so it is an isomorphism with the
prescribed labels.
Its inverse is the required section.
This exact injection also fixes the
isomorphism type of the finite lift;
no alternative central extension
is substituted for $Q$.

\emph{Uniqueness above each invariant class.}
The section just constructed identifies
$N=C\rtimes H$ and
\[
                   J=C\rtimes\langle j\rangle .        \tag{D.57j}
\]
The Hochschild--Serre sequence, with
$H^1(C,\F_2)=0$ and
$\operatorname{cd}_2(C)\leq1$, gives
the actual inflation isomorphism
\[
                   H^*(J,\F_2)\cong H^*(C_2,\F_2).
                                                               \tag{D.57k}
\]
The full Brown restriction theorem
\cite[Theorem~2 and \S1.1]{A-KochloukovaZalesskii2007}
applies: $C$ is open of finite
two-cohomological dimension and no
finite subgroup of $J$ can be a
four-group, since it injects into
$C_2$.
In high degree the source of its
dense restriction map is the finite
one-dimensional space in (D.57k).
Restriction to the displayed section
in (D.57j) is already nonzero in
every degree, so this map is injective.
Its image is therefore closed as
well as dense, so the restriction
is an isomorphism onto the product
of the section-normalizer cohomologies.
For a section $S\leq J$ its normalizer
is $C_C(S)\times S$, by the quotient
labels. Each such factor contributes
\[
                          1+\dim H^1(C_C(S),\F_2)
\]
in every sufficiently high degree.
There is consequently exactly one
conjugacy class of sections in $J$,
and its fixed group has zero first
cohomology over $\F_2$.
These are also exactly the
$C$-conjugacy classes: multiply a
given conjugator by the section
element with its quotient label.

Full sections of $E_Q$ containing
the chosen $T$ or $B$ lie in $N$,
and are exactly the full inverse
images of the order-two sections
of $J$. This follows from their
orders and the same kernel calculation
used above.
They therefore form one $C$-conjugacy
class. Every full four-group section
has a positive subgroup $T$ in one
of the invariant classes of
(D.57e); distinct such classes
cannot be identified by $F$-conjugacy.
The numbers of full section classes
are thus exactly three and one
as asserted. Every full dihedral
section has a rotation subgroup
in the unique class of D.53, so
the dihedral count is exactly one.

The fixed group $C_C(j)$ in the
Brown calculation is precisely
$C_F(H)$ in the original extension.
For example, commuting with $j$
modulo $T$ or $B$ means that a
commutator lies in that finite
subgroup. The commutator also lies
in $F$, so is trivial. The elements
of $C$ already commute with $T$
or $B$. This proves the cohomology
part of (D.57c).
Finally an element normalizing $H$
can be multiplied by the inverse
of the unique element of $H$ with
its outer label. The result lies
in $F$ and normalizes $H$, and
therefore centralizes each element
of $H$ by preservation of labels.
This gives the exact direct product
in (D.57c).

The program
\path{code/part_a/torus_two_group_lifts_d57.py}
checks all affine involutions of
$\F_2^2$ together with their invariant
three-point sets and verifies the
three-versus-one fixed-class count.
The full projectivity, Brown and
finite-preimage arguments above
provide the labelled lifts and
their uniqueness statements.
\end{proof}
